\documentclass[12pt]{article}
\usepackage[T1]{fontenc}
\usepackage[utf8]{inputenc}
\usepackage[brazil]{babel}
\usepackage[margin=1in]{geometry}
\setlength{\parindent}{0pt}
\begin{document}
\section*{Tema 82}
Processo: PEDILEF 2009.72.50.002288-2/ SC\\
Situacao: Julgado\\
Ramo: DIREITO ADMINISTRATIVO\\
Questao: Saber se a impugnação referente a ato de progressão/promoção de servidor público se sujeita ao prazo decadencial do art. 1° do Decreto n. 20.910/32.\\
Tese: A progressão/promoção de servidor público constitui ato único de efeito concreto, sendo o prazo decadencial para sua impugnação de 05 (cinco) anos (Decreto n. 20.910/32, art. 1º), contados a partir de sua publicação.\\
Fonte: http://pesquisa.in.gov.br/imprensa/jsp/visualiza/index.jsp?jornal=1\&pagina=146\&data=09/08/2012
\end{document}
